% festival_strategy.tex — Festival positioning

\section{Festival Strategy}

\subsection{Primary Target: Venice}

Venice is the ideal premiere. It is the festival most closely associated with prestige European art cinema, and its track record of launching Best Picture and Best Actor campaigns is unmatched: \textit{Nomadland}, \textit{Joker}, \textit{The Shape of Water}, \textit{T\'{a}r}, \textit{All Quiet on the Western Front}, \textit{Poor Things}. It is also the right contextual frame for this particular film: Central European setting, severe formal control, actor-forward seriousness, and enough emotional accessibility in the Berlin material to travel beyond purely austere cinephile response.

A Venice Competition selection provides:
\begin{itemize}[leftmargin=*]
\item European cultural credibility appropriate to the material
\item Golden Lion potential, which translates directly to Oscar momentum
\item Critical mass of awards-season press and industry attention
\item The ``Venice launch'' narrative that has become the most reliable path to Oscar viability
\item A premiere context where Brody's performance, the title object, and the controlled formal pressure will all read as assets rather than liabilities
\end{itemize}

\subsection{Alternative Primary: Cannes}

If the timeline or director attachment suggests Cannes over Venice, the film is equally suited to Competition there. The formal ambition, European sensibility, and director-driven packaging align with Cannes' programmatic identity. However, the May premiere date requires an earlier post-production schedule and a longer campaign runway to Oscar season, and Cannes can be slightly less forgiving to actor-forward prestige work if the directorial frame is not fully locked.

\subsection{Telluride / TIFF}

Regardless of Venice or Cannes premiere, the film should play Telluride and TIFF in the same festival cycle. Telluride provides the intimate, tastemaker-driven launch that converts festival buzz into awards-season velocity and is especially valuable for consolidating the Brody narrative. TIFF provides the industry-facing platform and the People's Choice Award pathway, and is useful proof that the film's emotional shape reaches beyond pure art-house severity.

\subsection{New York Film Festival}

A NYFF Main Slate selection would be a strong complement to the Venice/Telluride/TIFF circuit, reinforcing the film's literary prestige, formal credibility, and New York critical positioning in the domestic market.

\subsection{Post-Festival Release Strategy}

The optimal release pattern is a fall platform release following the festival circuit:

\begin{itemize}[leftmargin=*]
\item Venice/Telluride/TIFF (September)
\item NYFF (October)
\item Limited theatrical release in major markets (November)
\item Expanded theatrical and awards-qualifying run (December)
\item FYC screenings and campaign through February
\end{itemize}

\subsection{International Festival Positioning}

The film's Central European setting, literary pedigree, and co-production logic provide natural positioning at:
\begin{itemize}[leftmargin=*]
\item Karlovy Vary (post-premiere special screening; Czech connection)
\item Berlin (if premiere timing shifts; strong Mitteleuropa affinity)
\item London Film Festival (UK distribution launch)
\item European Film Awards consideration
\end{itemize}

\subsection{Sales Agent Positioning}

A European sales agent with strong festival relationships is essential. The Match Factory, Protagonist Pictures, and MK2 Films are the natural tier-one options, each with demonstrated expertise in prestige European drama, serious auteur packaging, and the kind of international positioning this film will require.
